\documentclass[12pt,a4paper]{article}
\usepackage[utf8]{inputenc}
\usepackage[T1]{fontenc}
\usepackage[spanish]{babel}
\usepackage{amsmath,amssymb,amsfonts}
\usepackage{geometry}
\usepackage{booktabs}
\usepackage{hyperref}
\geometry{margin=2.5cm}

\title{Una Perspectiva Causal--Horizonte sobre el Universo Temprano}
\author{Kevin Mu\~noz}
\date{}

\begin{document}
\maketitle

\begin{abstract}
Presentamos un marco conceptual centrado en el horizonte para interpretar la cosmología del universo temprano, en el cual el universo observable es tratado como un dominio causalmente acotado definido por una estructura de horizonte primordial. El marco no introduce nuevas leyes dinámicas ni modifica la Relatividad General. En cambio, reorganiza la interpretación de los fenómenos cosmológicos establecidos en términos de accesibilidad causal, condiciones de borde y entropía asociada al horizonte.

Dentro de esta perspectiva, la geometría del espacio-tiempo y la dinámica cosmológica de bulk se consideran descripciones efectivas restringidas por horizontes causales. Enfoques existentes sugieren que la evolución cosmológica estándar puede reformularse consistentemente en términos de termodinámica de horizontes, preservando la compatibilidad con la fenomenología convencional.

El fondo cósmico de microondas se interpreta tanto como un campo de radiación remanente asociado a la recombinación como una interfaz observacional correspondiente a una frontera causal. La expansión inflacionaria se trata igualmente como una descripción efectiva de regímenes en los que la evolución del horizonte domina el comportamiento de los grados de libertad accesibles.

No se asume ningún modelo microscópico específico. El marco identifica en cambio un conjunto mínimo de condiciones de consistencia estructural ---cierre causal, entropía asociada al borde y dinámica efectiva restringida por el horizonte--- que cualquier realización formal compatible debe satisfacer.
\end{abstract}

\setcounter{section}{-1}
\section{Posición Estructural del Marco}

El presente marco opera en tres niveles distintos pero compatibles:

\paragraph{Nivel ontológico:} Los dominios físicos se tratan como regiones causalmente auto-consistentes definidas por la accesibilidad causal y la estructura del horizonte.

\paragraph{Nivel interpretativo:} Los fenómenos cosmológicos observados se reconsideran como manifestaciones de condiciones de borde asociadas a horizontes causales.

\paragraph{Nivel dinámico efectivo:} La evolución cosmológica estándar puede admitir reformulaciones en términos de variables del horizonte y relaciones termodinámicas, sin modificación de la Relatividad General.

Estos niveles no son independientes. La capa ontológica restringe las interpretaciones admisibles, mientras que los supuestos interpretativos acotan el espacio de posibles realizaciones efectivas.

El presente trabajo se centra principalmente en el nivel interpretativo, permaneciendo estructuralmente compatible con formulaciones tanto ontológicas como efectivas.

\section{Introducción y Motivación}

La cosmología moderna proporciona una descripción cuantitativamente exitosa del universo observable dentro del marco $\Lambda$CDM. Combinado con una fase temprana de expansión acelerada, el modelo estándar da cuenta exitosamente del fondo cósmico de microondas (FCM), la formación de estructura a gran escala y la historia de expansión observada.

A pesar de este éxito empírico, varias preguntas fundamentales permanecen abiertas a nivel conceptual. Estas incluyen la interpretación de las condiciones iniciales, el estatus y origen de la inflación, el papel de la entropía en la evolución cosmológica y el significado de la estructura causal.

De manera independiente, los desarrollos en física gravitacional han establecido que los horizontes causales poseen propiedades termodinámicas, incluyendo temperatura y entropía proporcional al área del horizonte. Tales características sugieren que las fronteras causales pueden desempeñar un papel estructuralmente significativo en la determinación de la descripción física disponible para los observadores.

El presente trabajo no intenta resolver estos problemas mediante la introducción de nuevos mecanismos dinámicos. En cambio, explora si una reorganización de la interpretación cosmológica en torno a la accesibilidad causal y la estructura del horizonte puede proporcionar una perspectiva conceptual coherente, permaneciendo plenamente compatible con la fenomenología establecida.

El marco presentado aquí debe, por tanto, entenderse como conceptual e interpretativo más que predictivo o dinámico.

\section{Los Horizontes como Estructuras Físicas Fundamentales}

Dentro de la Relatividad General, un horizonte es una frontera causal, no un objeto material. Los horizontes delimitan regiones inaccesibles para clases específicas de observadores.

En múltiples contextos gravitacionales, los horizontes exhiben propiedades universales:

\begin{itemize}
  \item los horizontes están asociados con temperaturas efectivas
  \item los horizontes poseen entropía proporcional al área
  \item las descripciones físicas dependen de nociones de accesibilidad dependientes del observador
\end{itemize}

Estas propiedades surgen en espacio-tiempos de agujeros negros, contextos cosmológicos y marcos de observadores acelerados.

Desde una perspectiva centrada en el horizonte, tales resultados sugieren que las fronteras causales poseen significancia estructural más allá de su papel como características de coordenadas.

El presente marco adopta los horizontes causales como estructuras organizadoras físicamente significativas, sin comprometerse con una interpretación microscópica particular.

\section{Estatus y Alcance de las Afirmaciones}

El marco desarrollado aquí es de naturaleza conceptual e interpretativa.

No introduce:

\begin{itemize}
  \item nuevas ecuaciones dinámicas fundamentales
  \item grados de libertad físicos adicionales
  \item modificaciones de la Relatividad General
  \item predicciones observacionales independientes
\end{itemize}

Las afirmaciones realizadas en este trabajo deben, por tanto, interpretarse como condiciones de consistencia o reinterpretaciones de resultados existentes.

En particular:

\begin{itemize}
  \item no se realiza ningún supuesto respecto al origen microscópico de la entropía del horizonte
  \item no se propone ningún modelo específico de dinámica del horizonte primordial
  \item la fenomenología inflacionaria no es reemplazada
  \item la física estándar de recombinación permanece sin cambios
\end{itemize}

Las referencias a emergencia o derivación deben entenderse al nivel de descripción efectiva, no como afirmaciones de reducción fundamental.

\section{El Fondo Cósmico de Microondas como Frontera Causal}

Dentro de la cosmología estándar, el fondo cósmico de microondas se entiende como radiación remanente emitida durante la recombinación y posteriormente desplazada hacia el rojo a través de la expansión cosmológica.

El presente marco retiene esta descripción en su totalidad.

Sin embargo, al FCM se le puede asignar adicionalmente un papel interpretativo complementario.

Dentro de una perspectiva centrada en el horizonte, el FCM puede verse como poseedor de un estatus dual:

\begin{itemize}
  \item un campo de radiación asociado a la recombinación
  \item una interfaz observacional asociada a una frontera causal
\end{itemize}

Estos roles no son equivalentes y no deben identificarse.

El marco no afirma que el FCM sea generado por un horizonte, ni que la recombinación pueda ser reemplazada por física del horizonte. En cambio, asigna significancia estructural al hecho de que las observaciones presentes acceden al universo temprano a través de una superficie causalmente definida.

Desde este punto de vista, la uniformidad del FCM a gran escala puede interpretarse como compatible con condiciones de borde causales compartidas.

Aquí no se propone ningún mapeo cuantitativo entre las propiedades del horizonte y las anisotropías observadas. Establecer tal correspondencia sigue siendo un problema abierto.

\section{La Inflación como Fenómeno de Horizonte Emergente}

La inflación sigue siendo una descripción efectiva fenomenológicamente exitosa de la dinámica del universo temprano.

El presente marco no cuestiona la fenomenología inflacionaria ni intenta reemplazar los modelos existentes.

En cambio, la inflación puede interpretarse como la descripción de regímenes en los que la evolución de la accesibilidad causal y las cantidades relacionadas con el horizonte domina el comportamiento del dominio observable.

Esta interpretación no debe entenderse como la introducción de un mecanismo inflacionario alternativo.

Diferentes escenarios inflacionarios pueden corresponder, en cambio, a distintas realizaciones efectivas de estructuras similares centradas en el horizonte.

El marco, por tanto, reinterpreta el estatus conceptual de la inflación preservando su papel empírico.

En la realización efectiva, la variable de perturbación fundamental es $\psi = \delta R_A / R_A$, la fluctuación fraccional del radio del horizonte aparente. Para geometría espacialmente plana, las perturbaciones de densidad se relacionan con las fluctuaciones del horizonte mediante $\delta\rho/\rho = -2\psi$, proporcionando una descripción basada en el borde del origen de la estructura.

\section{Entropía, Información y la Flecha del Tiempo}

Las asociaciones entre entropía y horizontes están bien establecidas dentro de la termodinámica gravitacional.

Extendiendo estos conocimientos a contextos cosmológicos, el presente marco trata la entropía asociada al horizonte como una cantidad estructuralmente relevante para describir el dominio observable.

La flecha termodinámica del tiempo puede interpretarse como el reflejo de cambios en la accesibilidad causal y el crecimiento de la información disponible.

Esta afirmación debe entenderse como un reencuadre interpretativo, no como una derivación a partir de primeros principios.

No se afirma que esta perspectiva reemplace las descripciones termodinámicas estándar ni elimine la necesidad de condiciones cosmológicas especiales.

\section{Dimensionalidad y la Emergencia del Espacio}

Cerca de regímenes dominados por estructura del horizonte, la descripción espacial clásica puede dejar de ser fundamental.

Dentro de varios enfoques de gravedad cuántica, surgen indicaciones de reducción dimensional efectiva a energías suficientemente altas.

El presente marco permanece compatible con tales posibilidades y considera la emergencia de la geometría clásica como potencialmente asociada a la evolución de la estructura causal.

No se propone ningún mecanismo específico.

\section{Implicaciones Conceptuales y Expectativas de Consistencia}

El marco genera varias expectativas conceptuales:

\begin{itemize}
  \item la homogeneidad a gran escala puede verse como compatible con condiciones de borde compartidas
  \item el comportamiento inflacionario puede admitir interpretaciones centradas en el horizonte
  \item la entropía puede estar asociada principalmente a fronteras causales
  \item los dominios observables pueden caracterizarse naturalmente mediante accesibilidad causal
  \item la asimetría temporal puede emerger localmente en lugar de globalmente
  \item las realizaciones efectivas del marco producen un espectro de perturbaciones casi invariante de escala ($n_s \approx 1$) bajo supuestos estadísticos mínimos
\end{itemize}

Estas afirmaciones deben entenderse como requisitos de consistencia conceptual, no como predicciones observacionales.

\section{Problemas Abiertos y Direcciones Futuras}

Varias preguntas permanecen intencionalmente abiertas:

\begin{itemize}
  \item descripción microscópica de horizontes primordiales
  \item mapeo cuantitativo borde-a-bulk más allá del límite súper-horizonte
  \item origen de las escalas térmicas
  \item relación entre la evolución del horizonte y la inflación convencional
  \item determinación dinámica de $\varepsilon$ desde dentro del marco
  \item corrimiento del tilt: expresión termodinámica del horizonte para $\dot{\varepsilon}$
\end{itemize}

Los siguientes puntos listados como problemas abiertos en versiones anteriores del marco han sido resueltos en los documentos de formalización compañeros:

\begin{itemize}
  \item Formulación gauge-invariante de $\mathcal{R}$ en términos de $\psi$: derivada exactamente como $\mathcal{R} = [(1+\varepsilon)/\varepsilon]\,\psi$ en escalas súper-horizonte (\textit{Gauge-Invariant Formulation of the Horizon Perturbation}).
  \item Mecanismo de generación de perturbaciones: derivado a partir de la cuantización canónica del campo del horizonte con condiciones iniciales de Bunch--Davies (\textit{Scale Invariance from Horizon Field Quantization}).
  \item Sector tensorial: $n_T = -2\varepsilon$, $r = 16\varepsilon$, derivados a partir de los modos TT sobre el borde del horizonte (\textit{Tensor Perturbations from the Horizon Boundary}).
\end{itemize}

Estos problemas definen un programa de investigación futuro, no deficiencias del presente marco.

\section{Conclusión}

Hemos presentado un marco conceptual centrado en el horizonte para reorganizar la interpretación de la cosmología del universo temprano en torno a la estructura causal y las condiciones de borde.

El marco está intencionalmente limitado en su alcance. No introduce nueva dinámica ni propone un modelo cosmológico completo.

Su papel principal es identificar principios de consistencia estructural que alineen el razonamiento cosmológico con los conocimientos de la termodinámica de horizontes y la accesibilidad causal.

El trabajo futuro determinará si estas ideas admiten realizaciones cuantitativas capaces de reproducir fenómenos observacionales dentro de marcos predictivos.

\end{document}
